\documentclass{article}

% -----------------------------------------------------------------------------
% ICLR template
% Place the official `iclr2026_conference.sty` in the same directory.
% Official template: https://github.com/ICLR/Master-Template/raw/master/iclr2026.zip
% The fallback below is only to make this source compile when the style file is
% not yet present; remove the fallback for an actual ICLR submission.
% -----------------------------------------------------------------------------
\IfFileExists{iclr2026_conference.sty}{%
  \usepackage{iclr2026_conference,times}
}{%
  \usepackage[margin=1in]{geometry}
  \usepackage{times}
}

\usepackage{amsmath,amssymb,mathtools}
\usepackage{booktabs}
\usepackage{array}
\usepackage{multirow}
\usepackage{microtype}
\usepackage{algorithm}
\usepackage{algpseudocode}
\usepackage{enumitem}
\usepackage{xcolor}
\usepackage{hyperref}

\newcommand{\Df}{\mathcal{D}_{\mathrm f}}
\newcommand{\V}{\mathcal{V}}
\newcommand{\Mset}{\mathcal{M}}
\newcommand{\Aset}{\mathcal{A}}
\newcommand{\KL}{D_{\mathrm{KL}}}
\newcommand{\teacher}{\theta_0}
\newcommand{\student}{\theta}
\newcommand{\pteacher}{p_{\teacher}}
\newcommand{\pstudent}{p_{\student}}
\newcommand{\Qans}{Q_{\Aset}}

\title{Trajectory-Selective Distillation for Retain-Set-Free\\
Masked Diffusion Language Model Unlearning}

\author{Anonymous Authors}

\begin{document}
\maketitle

\begin{abstract}
We study retain-set-free machine unlearning for masked diffusion language models
(MDLMs). The central design principle is to treat a token prediction under a
particular mask configuration as the atomic denoising decision. We directly
modify the frozen teacher distribution only at answer-bearing denoising
positions, while self-distilling the original teacher at all remaining masked
positions. To account for the bidirectional generation dynamics of MDLMs, the
forgetting signal is applied across partially masked answer states rather than
only at a fully masked response. This document gives the notation and algorithm
for the resulting trajectory-selective distillation objective.
\end{abstract}

\section{Problem Setup}
Let $p_{\theta_0}$ be a pretrained or instruction-tuned MDLM and let
$\Df=\{(x,y,\Aset)\}$ be a forget set. Here, $x$ denotes the conditioning
prompt, $y=(y_1,\ldots,y_L)$ denotes the response, and
$\Aset\subseteq\{1,\ldots,L\}$ contains the answer-bearing token positions
whose information should be forgotten. We assume no retain examples are
available during unlearning.

An MDLM operates on partially masked sequences. For a mask set
$\Mset\subseteq\{1,\ldots,L\}$, let $y_{\Mset}$ denote the corrupted response
obtained by replacing $y_i$ with the mask token for every $i\in\Mset$. At a
masked position $i\in\Mset$, the model predicts
\begin{equation}
    p_{\theta}(\cdot\mid x,y_{\Mset},\Mset).
\end{equation}
The mask set is therefore part of the denoising state: the same answer token can
be predicted under many different amounts of visible left and right context.

\section{Notation}

\begin{table}[t]
\centering
\caption{Notation used by trajectory-selective distillation.}
\label{tab:notation}
\small
\setlength{\tabcolsep}{5pt}
\renewcommand{\arraystretch}{1.15}
\begin{tabular}{>{\raggedright\arraybackslash}p{0.19\linewidth}
                >{\raggedright\arraybackslash}p{0.73\linewidth}}
\toprule
\textbf{Symbol} & \textbf{Meaning} \\
\midrule
$\Df$ & Forget dataset containing prompt--response pairs and answer-bearing spans. \\
$(x,y)$ & Forget prompt $x$ and response $y=(y_1,\ldots,y_L)$. \\
$\Aset$ & Set of answer-bearing token positions in $y$. \\
$\Mset$ & Set of response positions masked in the current denoising state. \\
$y_{\Mset}$ & Partially masked response associated with mask set $\Mset$. \\
$\Qans(\Mset\mid\Aset)$ & Answer-stratified mask distribution, constrained so that $\Mset\cap\Aset\neq\varnothing$. \\
$\theta_0$ & Frozen original MDLM parameters; used as the self-distillation teacher. \\
$\theta$ & Trainable MDLM parameters, initialized from $\theta_0$. \\
$p^0_{\Mset,i}$ & Frozen-teacher distribution $p_{\theta_0}(\cdot\mid x,y_{\Mset},\Mset)$ at masked position $i$. \\
$\mathcal{B}_{\Mset,i}$ & Set of vocabulary tokens to demote at an answer-bearing decision. \\
$q_{\Mset,i}$ & Edited teacher target distribution at state $(\Mset,i)$. \\
$w_{\Mset,i}$ & Optional per-decision weight used to balance answer and non-answer masked positions. \\
$\V$ & Vocabulary. \\
\bottomrule
\end{tabular}
\end{table}

\section{Trajectory-Selective Distillation}

\subsection{Answer-stratified mask sampling}
Standard diffusion corruption may produce many states in which the answer is
fully visible and hence provides no direct forgetting signal. We instead sample
mask states from a distribution $\Qans$ satisfying
\begin{equation}
    \Mset\cap\Aset\neq\varnothing.
\end{equation}
For a multi-token answer, the sampler should cover different numbers of visible
answer tokens. A simple implementation first samples
\begin{equation}
    k \sim \operatorname{Uniform}\{1,\ldots,|\Aset|\},
\end{equation}
then masks $k$ uniformly chosen positions from $\Aset$. The remaining
non-answer positions are masked according to the backbone's ordinary diffusion
corruption distribution. This exposes the model to both fully masked and
partially revealed answer states.

\subsection{Selective teacher editing}
For every masked position $i\in\Mset$, compute the frozen-teacher distribution
\begin{equation}
    p^0_{\Mset,i}(v)
    =
    p_{\theta_0}(v\mid x,y_{\Mset},\Mset),
    \qquad v\in\V.
    \label{eq:teacher}
\end{equation}
We then construct a single target distribution $q_{\Mset,i}$ according to
whether the current denoising decision predicts an answer-bearing position.

For a non-answer position, we preserve the original model exactly:
\begin{equation}
    q_{\Mset,i}(v)=p^0_{\Mset,i}(v),
    \qquad i\notin\Aset.
    \label{eq:retain-target}
\end{equation}

For an answer-bearing position, let $\mathcal{B}_{\Mset,i}$ be the set of
vocabulary items to demote. The minimal version uses only the ground-truth
forgotten token,
\begin{equation}
    \mathcal{B}_{\Mset,i}=\{y_i\}.
    \label{eq:demote-set-minimal}
\end{equation}
The edited teacher removes the demoted probability mass and renormalizes the
remaining distribution:
\begin{equation}
q_{\Mset,i}(v)
=
\begin{cases}
0,
& v\in\mathcal{B}_{\Mset,i},\\[4pt]
\displaystyle
\frac{p^0_{\Mset,i}(v)}
{1-\sum_{u\in\mathcal{B}_{\Mset,i}}p^0_{\Mset,i}(u)},
& v\notin\mathcal{B}_{\Mset,i}.
\end{cases}
\label{eq:edited-target}
\end{equation}
A belief-demotion extension can enlarge $\mathcal{B}_{\Mset,i}$ to include a
small high-probability set under the frozen teacher, but this is not required by
the core method.

\subsection{Unified objective}
The trainable model is fitted to the selectively edited teacher using one KL
objective:
\begin{equation}
\boxed{
\mathcal{L}_{\mathrm{TSD}}(\theta)
=
\mathbb{E}_{(x,y,\Aset)\sim\Df}
\mathbb{E}_{\Mset\sim\Qans(\cdot\mid\Aset)}
\left[
\sum_{i\in\Mset}
    w_{\Mset,i}
    \KL\!\left(
        q_{\Mset,i}
        \,\|\,
        p_{\theta}(\cdot\mid x,y_{\Mset},\Mset)
    \right)
\right].
}
\label{eq:tsd}
\end{equation}
A simple balanced weighting scheme assigns half of the total weight to
answer-bearing masked decisions and half to non-answer masked decisions:
\begin{equation}
\sum_{i\in\Mset\cap\Aset}w_{\Mset,i}=\frac{1}{2},
\qquad
\sum_{i\in\Mset\setminus\Aset}w_{\Mset,i}=\frac{1}{2},
\label{eq:balanced-weight}
\end{equation}
whenever both groups are nonempty. If one group is empty, its mass is assigned
to the other group.

For $i\in\Aset$ and the minimal demotion set in
Eq.~\eqref{eq:demote-set-minimal}, $q_{\Mset,i}(y_i)=0$. Thus the derivative of
the KL loss with respect to the student's target-token logit
$z_{\theta,\Mset,i}(y_i)$ is
\begin{equation}
    \frac{\partial}{\partial z_{\theta,\Mset,i}(y_i)}
    \KL(q_{\Mset,i}\|p_{\theta})
    =
    p_{\theta}(y_i\mid x,y_{\Mset},\Mset) > 0,
    \label{eq:nonzero-gradient}
\end{equation}
whenever the model assigns nonzero probability to the forgotten token. Gradient
descent therefore directly reduces its logit at every sampled answer-bearing
denoising state.

\begin{algorithm}[t]
\caption{Trajectory-Selective Distillation (TSD)}
\label{alg:tsd}
\begin{algorithmic}[1]
\Require Frozen MDLM $p_{\theta_0}$; forget set $\Df$; trainable parameters
$\theta\leftarrow\theta_0$; answer-stratified mask sampler $\Qans$; learning
rate $\eta$.
\For{each optimization step}
    \State Sample a minibatch $\mathcal{B}\subset\Df$.
    \State Initialize minibatch loss $\mathcal{L}\leftarrow 0$.
    \ForAll{$(x,y,\Aset)\in\mathcal{B}$}
        \State Sample $\Mset\sim\Qans(\cdot\mid\Aset)$ such that
        $\Mset\cap\Aset\neq\varnothing$.
        \State Construct the partially masked response $y_{\Mset}$.
        \State Compute frozen-teacher distributions
        $p^0_{\Mset,i}\leftarrow
        p_{\theta_0}(\cdot\mid x,y_{\Mset},\Mset)$ for all $i\in\Mset$.
        \ForAll{$i\in\Mset$}
            \If{$i\in\Aset$}
                \State Set $\mathcal{B}_{\Mset,i}\leftarrow\{y_i\}$.
                \State Form $q_{\Mset,i}$ by removing
                $\mathcal{B}_{\Mset,i}$ from $p^0_{\Mset,i}$ and
                renormalizing as in Eq.~\eqref{eq:edited-target}.
            \Else
                \State Set $q_{\Mset,i}\leftarrow p^0_{\Mset,i}$.
            \EndIf
        \EndFor
        \State Compute student distributions
        $p_{\theta}(\cdot\mid x,y_{\Mset},\Mset)$.
        \State $\mathcal{L}\leftarrow\mathcal{L}
        +\sum_{i\in\Mset}w_{\Mset,i}
        \KL(q_{\Mset,i}\|p_{\theta}(\cdot\mid x,y_{\Mset},\Mset))$.
    \EndFor
    \State Update $\theta\leftarrow\theta-\eta\nabla_{\theta}\mathcal{L}$.
\EndFor
\State \Return unlearned parameters $\theta$.
\end{algorithmic}
\end{algorithm}

\section{Minimal Experimental Ablation}
The core mechanism can be isolated using four variants:
\begin{enumerate}[leftmargin=*]
    \item \textbf{Single-state demotion:} demote answer tokens only when the
    entire answer span is masked.
    \item \textbf{Pathwise demotion:} apply the same target demotion over fully
    and partially masked answer states.
    \item \textbf{Pathwise demotion + TSD:} additionally self-distill the
    frozen teacher at all non-answer masked positions using
    Eq.~\eqref{eq:retain-target}.
    \item \textbf{Belief demotion:} replace $\{y_i\}$ by a small
    teacher-defined high-probability demotion set to test whether broader
    belief suppression reduces paraphrase leakage.
\end{enumerate}
The first diagnostic is whether pathwise demotion suppresses exact answer
reconstruction under all levels of partial answer visibility. Only after this
condition is met should retain utility and semantic paraphrase leakage be used
to evaluate the preservation and belief-demotion components.

\end{document}
